\section{Complete families and the cubic obstruction}\label{sec:complete-family}

We turn to Theorem~\ref{thm:main}. The first step is to exclude the
cubic alternative from a complete family. We begin with the character
identities that record mutual unbiasedness and completeness.

\subsection{Character identities of a complete family}

Suppose that
\(B_1,\ldots,B_7\) are mutually unbiased orthonormal bases.
For \(r\ne s\), define their transition matrix by
\[
 H_{rs}(i,k)=\sqrt6\,\langle b_{r,i},b_{s,k}\rangle.
\]
Unbiasedness and orthonormality imply that \(H_{rs}\) is Hadamard and that
\(H_{sr}=H_{rs}^*\). With a reference basis \(r\) fixed, we abbreviate
\(H_{rs}\) and \(g_{H_{rs}}\) to \(H_s\) and \(g_s\).

\begin{lemma}\label{lem:basic}
For each fixed reference basis the following identities hold.
First,
\begin{equation}\label{eq:order-two}
 g_s(d)=0\quad(d\in D),
 \qquad
 \sum_{s\ne r}g_s(a)=0\quad(\ord(a)=2).
\end{equation}
For arbitrary charges \(a,b\) and each coordinate \(i\),
\begin{equation}\label{eq:same-matrix}
 \sum_{j=0}^5 g_s(a+e_i-e_j)g_s(b-e_i+e_j)=g_s(a+b).
\end{equation}
The same identity holds with both shifts reversed. For distinct external
bases \(s,t\ne r\),
\begin{equation}\label{eq:different-matrices}
 \sum_{d\in D}g_s(a+d)g_t(b-d)=0.
\end{equation}
\end{lemma}

The complete-set projector identity used in the proof is the
second-moment identity for a complex projective $2$-design; see
Zauner \cite[Satz~2.19]{Zauner1999} and
Klappenecker--R\"otteler \cite[Theorem~3]{KlappeneckerRotteler2005}.
We give the projection argument in our normalization. The character
identities are part of the Fourier approach
\cite{Matolcsi2012,MatolcsiRuzsaWeiner2013}.

\begin{proof}
The first assertion in \eqref{eq:order-two} is row orthogonality.
For the second, let
\(P_{s,k}=b_{s,k}b_{s,k}^*\), and let \(\mathcal D_s^0\) be the complex
space of traceless matrices diagonal in \(B_s\). If \(s\ne t\), then
\[
 \tr\bigl((P_{s,k}-I/6)(P_{t,l}-I/6)\bigr)=0.
\]
The spaces \(\mathcal D_s^0\) are therefore pairwise orthogonal for the
Hilbert--Schmidt inner product. Each has dimension five, and their
dimensions sum to \(35=6^2-1\), the dimension of the traceless matrix
space. Orthogonal projection consequently gives
\begin{equation}\label{eq:projector-identity}
 \tr(XY)=\sum_{s=1}^7\sum_{k=0}^5
       \tr(XP_{s,k})\tr(YP_{s,k})
 \qquad(\tr X=\tr Y=0).
\end{equation}
This is a complex bilinear identity: it follows by pairing the
complex-linear projection formula for \(X\) with \(Y\).

Any order-two charge is a sum \(d_1+d_2\) of two members of \(D\).
For \(s\ne r\), write \(h_k\) for the \(k\)th column of \(H_s\).
In basis \(B_r\), choose off-diagonal matrix units \(X,Y\) whose traces
against \(P_{s,k}\), for \(s\ne r\), are respectively
\(h_k^{d_1}/6\) and \(h_k^{d_2}/6\). For example,
\(d_1=e_i-e_j\) corresponds to \(X=E_{ji}\).
Since \(d_1+d_2\) has order two, \(d_2\ne-d_1\), so
\(\tr(XY)=0\). The reference basis contributes zero to
\eqref{eq:projector-identity}, and each external basis contributes
\(g_s(d_1+d_2)/6\). This proves the second identity in
\eqref{eq:order-two}.

For \eqref{eq:same-matrix}, expand the two character sums and write
\(x,y\) for their sampled columns. The left-hand side is
\[
 \frac1{36}\sum_{x,y}x^ay^b\frac{x_i}{y_i}
                 \sum_{j=0}^5\frac{y_j}{x_j}.
\]
The inner sum is \(6\) when \(x=y\) and zero otherwise, by column
orthogonality. The result is \(g_s(a+b)\). Conjugating the inner
orthogonality relation proves the version with reversed shifts.

Finally, if \(x\) is a column of \(H_s\) and \(y\) a column of \(H_t\),
their inner-product squared modulus is six. Hence
\[
 \sum_{d\in D}(x/y)^d
 =\left|\sum_{i=0}^5\frac{x_i}{y_i}\right|^2-6=0.
\]
Expanding the products of \(g_s\) and \(g_t\) proves
\eqref{eq:different-matrices}.
\end{proof}

The identities for one matrix hold for every Hadamard matrix.
The order-two sum in \eqref{eq:order-two} uses the completeness of the
seven-basis set. The later mixed identities retain this completeness
information through averaging over the six external bases.

\subsection{The cubic-pair obstruction}

Brierley and Weigert proved the stronger fact that the standard basis
and the basis defined by $T/\sqrt6$ cannot be extended by a third
mutually unbiased basis
\cite[Section~4.1 and Table~1]{BrierleyWeigert2009Constructing}.
The following proposition gives an exact Fourier derivation of the
complete-family exclusion needed here.

\begin{proposition}\label{prop:cubic-exclusion}
In a set of seven mutually unbiased bases in $\C^6$, no transition
Hadamard matrix $H$ can have $H^{\circ2}$ Hadamard.
\end{proposition}

To prove the proposition, suppose that such a pair occurs. By
Proposition~\ref{prop:cubic-classification}, rephasing and reordering
its basis vectors allows us to take its transition matrix to be $T$.
These operations preserve every mutual unbiasedness condition.
Use the first basis of this pair as reference, and call the other
five transition matrices $K_1,\ldots,K_5$. Define
\[
 f(b)=\frac15\sum_{s=1}^5g_{K_s}(b),\qquad G(b)=g_T(b).
\]
Lemma~\ref{lem:basic} gives $f(b)=G(b)=0$ at order one, and
Equation~\eqref{eq:order-two} gives
\begin{equation}\label{eq:cubic-order-two}
 f(b)=-\frac15G(b)\qquad(\ord(b)=2).
\end{equation}
Equation~\eqref{eq:different-matrices}, applied with its second charge
equal to $-c$ and then averaged over $K_1,\ldots,K_5$, gives
\begin{equation}\label{eq:cubic-cross}
 \sum_{d\in D}f(a+d)\overline{G(c+d)}=0
 \qquad\text{for all balanced }a,c.
\end{equation}
Let
\begin{equation}\label{eq:cubic-S}
 S(b)=\frac1{6!}\sum_{\pi\in S_6}
       \Ree\bigl(f(\pi b)\overline{G(\pi b)}\bigr).
\end{equation}
Because $f(-b)=\overline{f(b)}$ and likewise for $G$, this average
also equals the average over the distinct elements of the signed
permutation orbit of $b$.

Every entry of $T$ has cube one, so $G(3d)=1$ for $d\in D$.
The signed permutation orbit of $3(e_5-e_0)$ is precisely
$\{3d:d\in D\}$. If $z$ is a uniformly sampled column of a uniformly
sampled matrix among $K_1,\ldots,K_5$, then $\E z^b=f(b)$. Consequently
\begin{equation}\label{eq:cubic-positivity}
 0\le\E\left|\sum_{i=0}^5z_i^3\right|^2
 =6+\sum_{d\in D}f(3d)
 =6+30S\bigl(3(e_5-e_0)\bigr).
\end{equation}
Thus this orbit average must be at least $-1/5$. We will use
\eqref{eq:cubic-cross} to compute it and contradict this lower bound.
The next lemma concerns $T$ alone; it lets us calculate the coefficients
of those identities one permutation orbit at a time.

\subsection{Symmetries of the cubic matrix}

The row-permutation symmetry below also follows from the semilinear
automorphism group determined by Gillespie, \'O Cath\'ain, and Praeger
\cite[Proposition~2]{GillespieOCathainPraeger2018}. We give a direct
proof using the five-cycle classification.

\begin{lemma}\label{lem:cubic-symmetry}
Every permutation of the rows of $T$ can be completed to a symmetry
of $T$ by column permutations, row and column phase factors, and,
if necessary, complex conjugation. Consequently $|G(b)|$ is constant
on each orbit of balanced charges under permutations and negation.
\end{lemma}

\begin{proof}
A symmetry operation consists of row and column permutations, a choice
of whether to conjugate, and row and column phase factors that send $T$
to itself. Normalize the phase factor of output row zero to one.
The all-one output row and column then determine all the other phase
factors, removing the freedom to multiply the row factors by a common
phase and the column factors by its inverse.

Choose an initial row and column in $36$ ways, choose whether to
conjugate, and dephase at the chosen row and column.
Proposition~\ref{prop:cubic-classification} applies again. After matching
the additional ones on the diagonal, there are exactly $10$ orders of the remaining rows that,
together with the matching column orders, make the plus graph the
specified five-cycle: choose its initial vertex and its direction.
This constructs $36\cdot2\cdot10=720$ symmetry operations. Their row and
column permutations recover the chosen initial row and column as the
source indices sent to position zero; together with the conjugation
choice, their remaining orders recover the chosen cycle isomorphism.
Thus these are distinct normalized operations.

We show that their induced row permutations are distinct. Work with
exponents modulo three, and write $t_j=(t_{ij})_{i=0}^5$ for column $j$
of the exponent matrix. A symmetry fixing the row order acts on the
set $X=\{t_0,\ldots,t_5\}$ as
\[
 X\longmapsto sX+u,\qquad s\in\{1,-1\}.
\]
Indeed, because row zero has only zero exponents, the column phase
factors are common and may be absorbed into the row factors. The
remaining row factors have exponents modulo three, and the image of
the zero column shows that $u\in X$.

For $j\ne0$, column $t_j$ has exactly one zero in its noninitial
positions, at position $j$. Hence $-t_j\notin X$: its zero would force
it to equal $t_j$, which it does not. A nonzero translation preserving
$X$ would contain both $u$ and $2u=-u$, and is therefore impossible.
Negation with $u=0$ is impossible for the same reason.

Finally suppose $s=-1$ and $u=t_x\ne0$. Choose a neighbor $y$ of $x$
on the five-cycle, and let $z$ be their common non-neighbor. Directly
from \eqref{eq:cubic-matrix}, $t_x-t_y$ has its unique noninitial zero
at position $z$: at the other two positions outside $\{x,y,z\}$,
exactly one of $x,y$ is adjacent to that vertex, so the differences
are $\pm2$, hence nonzero modulo three. Were $t_x-t_y$ a column, it
would have to be $t_z$. But at
positions $x,y$ the former has entries $-1,1$, while $t_z$ has entries
$-1,-1$, a contradiction modulo three. Thus the only symmetry fixing
the row order is the identity, after the irrelevant common phase
normalization. Two of the $720$ operations with the same row
permutation would have such a quotient, so their row permutations
are all distinct. They exhaust the $720$ possible row permutations.

For a balanced charge, column phases cancel in the defining
character, while row phases multiply it by a unit complex number.
Conjugation conjugates the character. These observations prove the
last assertion, including negation since $G(-b)=\overline{G(b)}$.
\end{proof}

\subsection{Coefficients of the cross identities}

Let
\[
 a=(-1,-1,0,0,1,1),\qquad c=(0,0,0,0,-1,1),
\]
and let $\mathcal O$ be the set of distinct simultaneous permutations
$(A,C)$ of $(a,c)$.
Apply \eqref{eq:cubic-cross} to every $(A,C)\in\mathcal O$, multiply
by $\overline{G(A-C)}$, sum, and take real parts. To record the
coefficients explicitly, set
\[
 Q(b)=\sum_{\substack{(A,C)\in\mathcal O\\b-A\in D}}
                 G(A-C)G(b-A+C).
\]
The resulting equality is
\[
 \sum_b\Ree\bigl(f(b)\overline{Q(b)}\bigr)=0,
\]
where the sum is finite. For a complex number $z$, write
$z^{[1]}=z$ and $z^{[-1]}=\overline z$. Conjugacy of $f$ permits
symmetrizing the coefficient as
$R(b)=(Q(b)+\overline{Q(-b)})/2$: replacing $b$ by $-b$ in the finite
sum above and averaging leaves it unchanged. Explicitly,
\begin{equation}\label{eq:cubic-R}
 R(b)=\frac12
 \sum_{\substack{(A,C)\in\mathcal O,\ \epsilon\in\{-1,1\}\\
                  \epsilon b-A\in D}}
 \bigl[G(A-C)G(\epsilon b-A+C)\bigr]^{[\epsilon]}.
\end{equation}
We have therefore obtained the weighted cross identity
\begin{equation}\label{eq:cubic-weighted-cross}
 \sum_b\Ree\bigl(f(b)\overline{R(b)}\bigr)=0.
\end{equation}
Its coefficients depend only on $T$. We now calculate them.

The following table lists representatives $w_j$ of every balanced
charge of order at most three, up to permutation and negation. The
integer triples in its last two columns satisfy
\begin{equation}\label{eq:cubic-count-interpretation}
 6G(w_j)=\sum_{m=0}^2N_m\omega^m,
 \qquad
 72R(w_j)=\sum_{m=0}^2U_m\omega^m.
\end{equation}
\begin{table}[htbp]
\centering
\small
\begin{tabular}{@{}rlcc@{}}
\toprule
$j$&$w_j$&$(N_0,N_1,N_2)$&$(U_0,U_1,U_2)$\\
\midrule
0&$(-3,0,0,0,0,3)$&$(6,0,0)$&$(0,0,0)$\\
1&$(-3,0,0,0,1,2)$&$(2,2,2)$&$(0,0,0)$\\
2&$(-3,0,0,1,1,1)$&$(3,0,3)$&$(0,0,0)$\\
3&$(-2,-1,0,0,1,2)$&$(4,1,1)$&$(60,42,42)$\\
4&$(-2,-1,0,1,1,1)$&$(2,2,2)$&$(144,144,144)$\\
5&$(-2,0,0,0,0,2)$&$(2,2,2)$&$(0,0,0)$\\
6&$(-2,0,0,0,1,1)$&$(3,3,0)$&$(180,180,72)$\\
7&$(-1,-1,-1,1,1,1)$&$(3,3,0)$&$(540,540,216)$\\
8&$(-1,-1,0,0,1,1)$&$(4,1,1)$&$(480,336,336)$\\
9&$(-1,0,0,0,0,1)$&$(2,2,2)$&$(576,576,576)$\\
10&$(0,0,0,0,0,0)$&$(6,0,0)$&$(0,0,0)$\\
\bottomrule
\end{tabular}
\caption{Exact cube-root counts for $G$ and $R$.}
\label{tab:cubic-counts}
\end{table}

Here is an explicit summation rule reproducing
Table~\ref{tab:cubic-counts}. Put
\[
 v(b,k)=\sum_{i=0}^5 b_i t_{ik}\pmod3,
\]
with residues in $\{0,1,2\}$. For each input $b=w_j$, initialize two
separate arrays $N,U$ of three zero integers and apply these loops:
\begin{quote}
\begin{verbatim}
for k = 0,...,5:
    N[v(b,k)] += 1
for p,q = 0,...,5 with p != q:
    for 0 <= h < i <= 5 with {h,i} disjoint from {p,q}:
        A = e_p + e_q - e_h - e_i
        u = 2*e_p - e_h - e_i
        for s in {-1,1} with s*b - A in D:
            for k,l = 0,...,5:
                U[(s*(v(u,k) + v(s*b-u,l))) mod 3] += 1
\end{verbatim}
\end{quote}
Indeed, the $30\binom42=180$ pairs in $\mathcal O$ are uniquely
parametrized by
\[
 A=e_p+e_q-e_h-e_i,\qquad C=e_q-e_p,
\]
with precisely the conditions used in the loops; thus $u=A-C$.
Each product of two $G$ terms contributes $36$ cube roots with
coefficient $1/36$, and the factor $1/2$ in
\eqref{eq:cubic-R} accounts for the denominator $72$ in
\eqref{eq:cubic-count-interpretation}.

For example, row $3$ gives
$G(w_3)=(4+\omega+\omega^2)/6=1/2$ and
$R(w_3)=(60+42\omega+42\omega^2)/72=1/4$.
Using $1+\omega+\omega^2=0$ in every row gives
\begin{equation}\label{eq:cubic-R-multipliers}
 R(b)=\lambda_jG(b),\qquad
 \lambda_3=\tfrac12,\quad \lambda_6=3,\quad
 \lambda_7=9,\quad \lambda_8=4,
\end{equation}
on the signed permutation orbit of $w_j$, and $R(b)=0$ on all other
orbits. To justify extending the representative calculation, observe
that every factor in \eqref{eq:cubic-R} has a balanced charge. Its
column phases therefore cancel. The two charges in each product add
to $\epsilon b$, so after the prescribed conjugation its row phase
factor is exactly the factor for $G(b)$. Row permutations preserve
the full set $\mathcal O$, and conjugation conjugates both $R$ and $G$.
Lemma~\ref{lem:cubic-symmetry} therefore gives the asserted covariance
on every orbit; also $R(-b)=\overline{R(b)}$. No other orders can occur
in the support of $R$, since $\epsilon b=A+d$ with
$\ord(A)=2$ and $d\in D$ implies $\ord(b)\le3$.

\subsection{The contradiction}

\begin{proof}[Completion of the proof of Proposition~\ref{prop:cubic-exclusion}]
Write $S_j=S(w_j)$ for the representatives in Table~\ref{tab:cubic-counts}.
The orbit sizes for $j=3,6,7,8$ are respectively
$360,120,20,90$. By \eqref{eq:cubic-R-multipliers}, their coefficients in
\eqref{eq:cubic-weighted-cross}
are $180,360,180,360$. Dividing by $180$ proves
\begin{equation}\label{eq:cubic-first-relation}
 S_3+2S_6+S_7+2S_8=0.
\end{equation}

Table~\ref{tab:cubic-counts} and Lemma~\ref{lem:cubic-symmetry} imply
\[
 S_1=S_4=S_5=S_9=0.
\]
At types $6$ and $8$ the squared modulus of $G$ is $1/4$, so
\eqref{eq:cubic-order-two} gives
\begin{equation}\label{eq:cubic-known-S}
 S_6=S_8=-\frac1{20}.
\end{equation}

Next take $a=c=w_5$ in \eqref{eq:cubic-cross} and average over
coordinate permutations. Do the same with $a=c=w_8$. Classifying the
$30$ charges $w_j+d$, $d\in D$, by their signed permutation orbits
gives the following counts:
\[
\begin{array}{c|rrrrrrrr}
 &0&1&3&4&6&7&8&9\\\hline
 w_5+D&1&8&12&0&8&0&0&1\\
 w_8+D&0&0&4&8&4&2&8&4.
\end{array}
\]
These counts can be obtained simply by adding each of the $30$
vectors $e_i-e_j$ to the displayed representatives and sorting the
coordinates, allowing a simultaneous change of sign. Averaging
over permutations turns each resulting charge into its orbit
average in \eqref{eq:cubic-S}. Substituting the known values gives
\begin{equation}\label{eq:cubic-last-relations}
 S_0+12S_3=\frac25,
 \qquad 4S_3+2S_7=\frac35.
\end{equation}
Together, \eqref{eq:cubic-first-relation},
\eqref{eq:cubic-known-S}, and \eqref{eq:cubic-last-relations} yield
\[
 S_3=S_7=\frac1{10},\qquad S_0=-\frac45.
\]

Since $w_0=3(e_5-e_0)$, equation~\eqref{eq:cubic-positivity} now gives
\[
 \E\left|\sum_{i=0}^5z_i^3\right|^2
 =6+30S_0=-18,
\]
contradicting nonnegativity. This proves the proposition.
\end{proof}
